\documentclass[11pt]{article}
\usepackage[a4paper,margin=27mm]{geometry}
\usepackage{amsmath,amssymb,amsthm,booktabs,hyperref}
\newtheorem{theorem}{Theorem}
\newtheorem{lemma}[theorem]{Lemma}
\newtheorem{corollary}[theorem]{Corollary}
\newcommand{\F}{\mathbb F}
\newcommand{\PSL}{\operatorname{PSL}}
\newcommand{\SL}{\operatorname{SL}}
\newcommand{\tr}{\operatorname{tr}}
\newcommand{\Span}{\operatorname{span}}
\title{Elations and two translation modules}
\author{A recovery of the even-characteristic classification in L027}
\date{}
\begin{document}\maketitle
Let $F$ be a finite field of characteristic two. A string C-marking is an
ordered independent generating tuple of involutions satisfying the string
commutation relations and the subgroup intersection property. It is
nondegenerate when every adjacent product has order greater than two.
We prove the assertions of Leemans--Vandenschrick, Section~5.4--5.5, with
the actual field and intersection conditions made explicit. The matrices
in the normal forms are those of that source. The common geometric
argument and the two-module proof below replace its separate calculations.

We retain the elementary projective interpretation of a linear action,
the usual criterion which deduces the intersection property from the two
end markings and their intersection, and Dickson's subgroup theorem in
the following precise form. A subgroup of $\PSL_2(F)$ is cyclic,
dihedral of order $2k$ with $k$ odd, contained in a point stabilizer, or
conjugate to a natural $\PSL_2(K)$ for a subfield $K\subseteq F$.
These alternatives include the small groups; in particular $A_5$ is
$\PSL_2(4)$, and an $A_4$ subgroup is in a point stabilizer. We also
retain the standard faithful action of the stabilizer of a nonsingular
plane conic as $\PSL_2(F)$. No subgroup multiplicity or source census
is used.

\section*{Commutation and the flag case}
Every projective involution has a unique order-two lift to $\SL_3(F)$,
since the scalar centre has odd order dividing three. The lift has the
form $I+u f$, where $u\ne0$, $f\ne0$ and $f(u)=0$. Its centre is the
point $\langle u\rangle$ and its axis is the line $\ker f$.
If two projective involutions commute, their order-two lifts commute:
an equality $AB=\zeta BA$ with scalar $\zeta$ implies $\zeta^2=1$.

\begin{lemma}\label{commute}
Two distinct elations commute if and only if they have the same centre
or the same axis.
\end{lemma}
\begin{proof}
For $N=u f$ and $M=v g$, commutation is
$u f(v)g=v g(u)f$. If both the vectors and the functionals are independent,
this forces $f(v)=g(u)=0$. The independent vectors $u,v$ would then lie in
the one-dimensional intersection of the distinct kernels, a contradiction.
Conversely, a common centre or axis makes both products zero.
\end{proof}

\begin{theorem}\label{triangular}
A nondegenerate string C-marking whose group fixes a point and an
incident line has rank at most three. The possibilities in ranks one,
two and three are respectively $C_2$, $D_8$, and $C_2^2\wr C_2$, of
types $[]$, $[4]$, and $[4,4]$.
\end{theorem}
\begin{proof}
Choose a basis along the fixed flag. Every generator is upper unitriangular.
Write such a matrix as $(x,y,z)=I+xE_{12}+yE_{13}+zE_{23}$; it is an
involution precisely when $xz=0$. A generator with $x=z=0$ is central and
cannot occur in a nondegenerate marking of rank at least two. Noncommuting
adjacent generators must alternate between $x\ne0,z=0$ and $x=0,z\ne0$.
Their product has order four, giving $D_8$. At rank four the first and
fourth generators are of opposite kinds and do not commute, contrary to
the string relation.

At rank three, reverse the coordinates if necessary and write the ends
as $r=(a,b,0)$, $t=(e,f,0)$, and the middle as $s=(0,c,d)$, with
$a,e,d\ne0$. Their commutators with $s$ are the central matrices
$z_a=(0,ad,0)$ and $z_e=(0,ed,0)$. If $a=e$, the two dihedral ends
share the nonidentity $z_a=z_e$, which is not in $\langle s\rangle$;
the intersection property fails. If $a\ne e$, the four commuting
involutions $r,t,z_a,z_e$ are independent over $\F_2$: independence of
$a,e$ first kills the coefficients of $r,t$, and then that of $ad,ed$
kills the remaining coefficients. They form a normal group of order
$16$. Conjugation by $s$ sends $r$ to $rz_a$ and $t$ to $tz_e$.
It therefore interchanges the two bases $(r,t)$ and $(rz_a,tz_e)$,
giving $C_2^2\wr C_2$. The two dihedral end subgroups meet in
$\langle s\rangle$, since their respective base spaces
$\langle r,z_a\rangle$ and $\langle t,z_e\rangle$ are disjoint.
This proves the intersection property and the asserted type.
\end{proof}

\section*{Rank three without an incident fixed flag}
We write $M(x,y)=I+xE_{21}+yE_{23}$. The following small module
calculation will be used twice. If $t\ne0$, $K=\F_2(t)$, and the two
involutions act on rows by
\[
 (x,y)\longmapsto(x,y+t x),\qquad
 (x,y)\longmapsto(x+y,y),
\]
the module generated by $(0,c)$ is $cK^2$. Indeed the differences
give $(c,0)$ and $(0,tc)$; iteration gives every polynomial in $t$
in both coordinates. The reverse inclusion follows from invariance.
The product of the two matrices has trace $t$ and determinant one.
Its eigenvalues are $\lambda,\lambda^{-1}$, with $\lambda\ne1$;
its order $k>1$ is odd, and $K=\F_2(\lambda+\lambda^{-1})$.
Consequently $k$ divides $|K|-1$ or $|K|+1$. The generated group is
the actual dihedral group $D_{2k}$.

\begin{theorem}\label{three}
Let $(r,s,t)$ be a rank-three string C-marking with no fixed incident
point--line pair. Up to a projective change of coordinates and duality,
its lifts have the form
\[
 r=I+aE_{13},\qquad s=I+E_{31},\qquad
 t=I+bE_{13}+cE_{23},\qquad a\ne0.
\]
These lifts fix $e_2$, so their group meets the scalar centre trivially.
If $b=0$ or $b=a$, then $c\ne0$ and the group is
$K^2\rtimes D_{2k}$, where $K=\F_2(a)$ and the action is the natural
two-dimensional dihedral action. The types are respectively $[k,4]$
and $[k,k]$. Otherwise the group is $\PSL_2(K)$ for
$K=\F_2(a,b)$.
\end{theorem}
\begin{proof}
The ends commute, so after duality they share an axis $l$. The axis $m$
of $s$ differs from $l$, since otherwise all three commute and fix a
flag. Put $p=l\cap m$. The centre of $s$ is not $p$: otherwise the
whole group preserves $(p,l)$. The centres of the two ends cannot both
be $p$, since then $(p,m)$ is fixed. Reversing the marking if necessary,
take the centres of $r,s$ and the point $p$ as the three coordinate
points. Scaling gives the displayed matrices.

If $b=0$, the third generator is $M(0,c)$; if $b=a$, its product with
$r$ is $M(0,c)$. In either case $c\ne0$ by nonidentity and independence.
The group generated by $r,s$ has trivial intersection with the translation
group $M(F^2)$, and its action on that group is precisely the preceding
module calculation. Thus the asserted semidirect product is an equality
of actual matrix groups. When $b=0$, the other dihedral end has order
eight and intersects $\langle r,s\rangle$ in $\langle s\rangle$.
When $b=a$, set $v=M(c/a,0)$. Conjugation by $v$ fixes $s$ and takes
$r$ to $t$, so the other end is the conjugate of $D=\langle r,s\rangle$.
Its intersection with $D$ is the stabilizer in $D$ of the nonzero vector
$(c/a,0)$. No nonidentity rotation fixes such a vector, and two distinct
reflections fixing it would have a nonidentity rotation as their product.
The stabilizer is therefore exactly $\langle s\rangle$. These arguments
also prove the claimed types and the intersection property for these
normal forms.

Now suppose $b\notin\{0,a\}$. If $c=0$, the action on the plane
$\langle e_1,e_3\rangle$ is faithful. If $c\ne0$, direct substitution
shows that all three matrices preserve the nonsingular conic
\[
 Q=\frac{x_1^2+a x_1x_3+a x_3^2}{b(a+b)}+\frac{x_2^2}{c^2}=0.
\]
Its polar radical is $\langle e_2\rangle$ and $Q(e_2)\ne0$.
The induced action on the quotient by this point is faithful: a kernel
element has the form $M(x,y)$, and preservation of $Q$ gives
$(x x_1+y x_3)^2/c^2=0$, hence $x=y=0$. In either case the group is
isomorphic to
\[
 H=\left\langle
 \begin{pmatrix}1&a\\0&1\end{pmatrix},
 \begin{pmatrix}1&0\\1&1\end{pmatrix},
 \begin{pmatrix}1&b\\0&1\end{pmatrix}\right\rangle
 \le\SL_2(K),\qquad K=\F_2(a,b).
\]
The first and third matrices are distinct commuting involutions, excluding
the odd dihedral alternative. The first and second have different unique
fixed points, excluding a point stabilizer. Dickson's theorem therefore
gives a conjugate of $\PSL_2(K')$. The traces of the products of the
first and second and of the second and third matrices are $a,b$.
Conjugacy preserves trace, so $K\subseteq K'$; the group order bound
inside $\SL_2(K)$ gives $|K'|\le|K|$. Thus $K'=K$.
The conclusion is a classification of groups of the stated C-markings;
it does not assert that every remaining pair $a,b$ has the intersection
property.
\end{proof}

\section*{Rank four and the actual intersection}
\begin{theorem}\label{four}
A nondegenerate rank-four string C-marking in $\PSL_3(F)$ is, up to
coordinates and duality, of the form
\[
 r_0=I+E_{23},\quad r_1=I+aE_{31},\quad
 r_2=I+bE_{13},\quad r_3=I+E_{21},\qquad ab\ne0.
\]
Put $K=\F_2(ab)$ and let $k$ be the order of $r_1r_2$.
These matrices satisfy the intersection property exactly when $a\notin K$,
equivalently $b\notin K$. In that case their group is
$K^4\rtimes D_{2k}$, their two ends are $K^2\rtimes D_{2k}$,
and their type is $[4,k,4]$. There is no nondegenerate marking of
rank greater than four.
\end{theorem}
\begin{proof}
The group cannot fix an incident flag by Theorem~\ref{triangular}.
After duality, $r_0,r_2$ share an axis $l$. If $r_1,r_3$ share a
centre $p$ on $l$, the whole group fixes $(p,l)$; if their common
centre is off $l$, then $r_0,r_3$ have neither a common axis nor a
common centre and cannot commute. Thus $r_1,r_3$ share an axis
$m\ne l$. Since $r_0,r_3$ commute, their common centre is
$p=l\cap m$. Nondegeneracy forces the centres of $r_1,r_2$ to
differ from $p$. Taking these three centres as coordinate points and
scaling gives the displayed form.

Let $D=\langle r_1,r_2\rangle$. On the coordinates $(x_1,x_3)$
it is the dihedral group of order $2k$, with rotation trace $ab$;
it meets $M(F^2)$ trivially. On this translation group conjugation is
\[
 (x,y)^{r_1}=(x+a y,y),\qquad
 (x,y)^{r_2}=(x,y+b x).
\]
The module generated by $r_0=M(0,1)$ is
$N_0=M(aK,K)$, and that generated by $r_3=M(1,0)$ is
$N_3=M(K,bK)$. To see equality, successive differences first give
the other coordinate and then multiplication by $ab$; iteration gives
all polynomials in $ab$. The displayed spaces are invariant, proving
both inclusions. Hence the two actual end groups are $N_0\rtimes D$
and $N_3\rtimes D$.

If $a\in K$, then $b\in K$ and $N_0=N_3$, so the end intersection
is larger than $D$ and IP fails. Otherwise $aK\cap K=0$ and
$bK\cap K=0$, giving $N_0\cap N_3=0$ and end intersection exactly
$D$. Each end has rank-three IP: for example
$\langle r_0,r_1\rangle$ has translation space
$\Span_{\F_2}\{(0,1),(a,0)\}$ and meets $D$ in
$\langle r_1\rangle$. Its outer adjacent order is four, and the
other end follows by reversal. The end criterion proves full IP.
The sum $N_0\oplus N_3$ has dimension four over $K$, which gives
the group and type claimed.

Finally a nondegenerate marking of rank at least five has a rank-four
initial segment of this form. Its fifth generator commutes with $r_0$
and $r_2$, which have a common axis and distinct centres. By
Lemma~\ref{commute}, it must have that same axis. It also commutes
with $r_1$, whose axis is different and whose centre is not on the
common axis. This is impossible by the same lemma.
\end{proof}

\section*{Scope and attribution}
Theorems~\ref{triangular}, \ref{three}, and \ref{four} recover the
source's Theorems~5.4, 5.5, and 5.6, and the corresponding rank-three
and rank-four rows of Table~2. In particular the notation for the
elementary abelian groups refers to the exact trace fields displayed
above; it does not identify an arbitrary divisor of the ambient field
order with a subfield. The classification of rank-three projective-line
C-markings remains its separately attributed earlier result. The present
argument classifies their possible group, without repeating that census.
The rank-two semisimple rows of Table~2 and the odd-characteristic
classifications are outside this component. No chiral rank-six
nonexistence result is asserted here: that argument still has to join
the actual two regular end actions.
\end{document}
